\documentclass[a4paper, 12pt]{article}

\usepackage[utf8]{inputenc}
\usepackage[T1]{fontenc}
\usepackage{textcomp}
\usepackage{amssymb}
\usepackage{newtxtext} \usepackage{newtxmath}
\usepackage{amsmath, amssymb}
\newtheorem{problem}{Problem}
\newtheorem{example}{Example}
\newtheorem{lemma}{Lemma}
\newtheorem{theorem}{Theorem}
\newtheorem{problem}{Problem}
\newtheorem{example}{Example} \newtheorem{definition}{Definition}
\newtheorem{lemma}{Lemma}
\newtheorem{theorem}{Theorem}
\DeclareMathAlphabet{\mathcal}{OMS}{cmsy}{m}{n}

\begin{document}

En la última página de mi diario, esto está escrito:

\begin{quote}
La noche anterior todavía circulaba en nuestros pensamientos como un río
caudaloso y taciturno. Mi silencio disimulaba la eterna tristeza de mi vida. El
suyo, no sé qué... Entonces expresó este sentimiento: «leeré para ti cuando tus
ojos ya no vean». (Las almas sensibles producen tan exquisita belleza...)
Es una de las cosas más hermosas que me han dicho.
\end{quote}

Una de las cosas que me encantan de vos es tu capacidad de querer, de producir
pensamientos de ese tipo espontáneamente, como si fueran ocurrencias comunes. 

~

Si me conocieras por años, lentamente, como que se conocen los misterios, tal
vez podrías ver la tristeza de mi corazón. Me paso la vida añorando encontrar
almas afines, almas que conocer desnudamente, con la diáfana transparencia de un
agua espiritual. En todas encuentro un charco pantanoso que no deja pasar la luz
de la mirada.

~ 

Es cierto que, a veces, las encontré. Una es mi mejor amigo. Una ha muerto para
siempre. Una está conmigo. Tres en veintisiete años no es decir mucho. Por eso,
desde que sentí que te conozco desde siempre, y empecé a sospechar que sos una
de ellas, quise tenerte conmigo. \textit{Eso} ---transparencia, sensibilidad y
luz--- es lo que busco en una relación. Y vos lo tenés en abundancia.

~ 

Yo, por otra parte, tal vez no pueda darte lo que buscás o necesitás. Si eso es
así, es tu deber abandonarme. Voy a ser el primero en celebrar que persigas
\textit{tu} camino, acorde a \textit{tus} convicciones, deseos y necesidades. No
quisiera conducirte ---más bien verte caminar tu propia senda, como un
respetuoso testigo de tu vida---.

~ 

Pero mientras en tu corazón exista una voz, o un eco de una voz, que diga mi
nombre, tendrás mi cariño. Incluso si desaparezco de tu vida. No se trata de si
quiero o decido dártelo: ya es tuyo. Te di algo, algo secreto e inconfesable, de
mi corazón. No es mucho, pero es verdadero. Y esas ofrendas son eternas. Vienen
sin devolución.

~
~ 

No sé cuál es el sentido de esta carta. No tiene un mensaje definido. Tal vez
es una forma de decir: \textit{Rojhayjhueteiva}. (El castellano no es suficiente.)

~ 

Nos deseo lo mejor. En cualquier caso, seguiré abrazándote y besándote a través
del tiempo. \textit{Rohechaga'u}.

~

---SLP



























\end{document}



